\newpage
\section{作文初稿：先看实际完成了什么}
\textbf{正文 166 词，标题不计入，符合 120--180 词要求。}统计时将 \en{college.So} 视为两个词，带撇号的形式按一个词计。本人表示写作用时不长，但未记录起止时间，不据截图时间推定具体耗时。

初稿回应了学生会征求建议的任务，有开头、主体和结尾。主体能辨认出三项想法：组织运动，帮助学生了解专业与未来工作的联系，以及了解学生的日常困难。主要问题是基础语法、搭配和长句组织；不是没有内容，也不是需要再凑字数。

\subsection{初稿存档}
以下按截图转录，去掉屏幕自动换行，保留可辨认的拼写、大小写、标点和语法错误；不是修改稿。
\begin{quote}\small
Some suggestion to enrich students life

I headed that stuent union is collecting suggention to improve work for the coming year, as a students of our college, I want to make some suggestion and make our school better in next year.

I believe the most important thing about next year is how to make students know why we need to study in college.So I think we can prepare some activities such as running around school and know that we can make our body strong before walk into our future work, or make an activity and let students to know how there major can help to get value In future. I believe that we can also prepare to walk into some students’s life to know how the students experience their life and what they to in a whole day, it can let student union know students’s life and organize more activities to teach students solve problems when they face problems in life.

I hope these suggestion can help students get some ideas.
\end{quote}

\subsection{先处理三类问题}
\begin{enumerate}
\item \textbf{句型准确。}重点是 \en{let somebody do}、\en{teach somebody to do} 和 \en{before doing}；不同动词的结构不能混用。
\item \textbf{一句表达一层意思。}开头和主体末尾都用逗号连接了独立句。建议、作用、下一项建议，可以分别成句。
\item \textbf{回到学生生活。}“明白为什么上大学”范围很大；需要落到学生会能组织的活动，以及学生能得到的实际帮助。
\end{enumerate}

\textbf{检查顺序：}内容与字数 $\rightarrow$ 句子是否完整 $\rightarrow$ 动词后面的结构 $\rightarrow$ 名词单复数、代词和拼写。写得快以后，把剩余时间用于这一轮检查。

\newpage
\section{逐句修改：从自己的话出发}
\subsection{标题和开头}
原题目中的 \en{Some suggestion to enrich students life} 可改为：
\en{Some Suggestions for Enriching Student Life}。这里有多项建议，\en{suggestions} 用复数；\en{student life} 可以整体表示学生生活。

原开头有两项独立表达：“听说学生会征求建议”和“作为学生提出建议”。拆开后是：
\begin{quote}
I have heard that the student union is collecting suggestions for the coming year. As a student at our university, I would like to offer some ideas to enrich student life.
\end{quote}
\begin{itemize}
\item \en{headed} 是 \en{head} 的过去式，不表示“听说”；这里用 \en{heard}，也可用 \en{I hear that ...}。
\item \en{stuent} 改为 \en{student}；\en{suggention} 改为 \en{suggestions}。
\item \en{as a students} 改为 \en{as a student}；\en{some suggestion} 改为 \en{some suggestions}。
\item \en{next year} 作时间状语时，前面通常不加 \en{in}。本稿开头用 \en{for the coming year}，不重复时间。
\end{itemize}

\subsection{主体第 1 句：把大目标改成可执行的建议}
\en{make students know why we need to study in college} 的意思可以理解，但它把题目扩大为“上大学的意义”。修改稿删去这一笼统总起句，让后文三项建议直接承担主体。专业与学习意义的想法保留在第 2 项建议中。

\subsection{主体第 2 句：拆开运动与专业活动}
原句从 \en{So I think ...} 一直写到 \en{In future}，把运动、健康、就业、专业价值放在了一起。先分成两项，再分别写活动和作用：
\begin{itemize}
\item 运动：\en{organize regular running activities on campus}；
作用：\en{help students stay healthy and make new friends}。
\item 专业：\en{invite graduates to share how they use what they learned in their majors at work}；
作用：帮助理解所学的价值，为未来职业作准备。
\end{itemize}
\textbf{语法修改：}\en{let students to know} 应为 \en{let students know}，因为 \en{let + 人 + 动词原形}；\en{there major} 中的 \en{there} 应为 \en{their}。\en{before walk} 若保留原结构，应改为 \en{before walking}。但“走进未来工作”整体也不自然，宜直接说 \en{prepare for their future careers}。

\textbf{内容展开：}“regular”“make new friends”和毕业生分享的具体形式，是根据原意补出的细节，不是对原句逐字纠错。\en{prepare some activities} 并非在所有语境中都错；这里表示“组织活动”，用 \en{organize} 更直接。

\subsection{主体第 3 句及结尾：明确调查的用途}
\en{walk into some students's life} 像是把“走进学生生活”直接翻译过去；可写为 \en{learn about students' daily lives and difficulties}。\en{students} 已以 s 结尾，所有格写成 \en{students'}，不是 \en{students's}。

\en{what they to} 改为 \en{what they do}；\en{teach students solve problems} 应为 \en{teach students to solve problems}。修改稿使用 \en{help students solve everyday problems}，意思相近，结构是 \en{help + 人 + (to) do}。再把“了解生活”具体化成调查，并把逗号连接的两个独立句拆开。

结尾的 \en{these suggestion} 改为 \en{these suggestions}；\en{help students get some ideas} 比较模糊，改为说明建议如何帮助学生会了解需求、丰富校园生活。

\newpage
\section{保留原意后的考场版本}
\textbf{正文 136 词，无需依靠标题补足字数。}三项想法来自初稿，活动形式及部分作用作了适度展开。本稿是修改示例，不是本人原始表达，也不作为官方范文或评分结果。

\begin{quote}
I have heard that the student union is collecting suggestions for the coming year. As a student at our university, I would like to offer some ideas to enrich student life.

First, the student union could organize regular running activities on campus. These activities would help students stay healthy and make new friends. Second, it could invite graduates to share how they use what they learned in their majors at work. This would help students understand the value of their studies and prepare for their future careers. Finally, the student union could conduct surveys to learn about students' daily lives and difficulties. Based on the results, it could organize practical activities to help students solve everyday problems.

I hope these suggestions will help the student union better understand students' needs and make campus life more enjoyable.
\end{quote}

\section{可灵活替换的结构}
模板保留的是句子的功能与语法。方括号表示需要按题意填入的部分，不能原样写进作文。

\begin{enumerate}
\item \textbf{说明身份与目的：}\\
\en{As a [身份], I would like to offer some suggestions on [名词或动名词短语].}\\
适用于确有相应身份的建议类回应；本题可填 \en{student} 与 \en{improving campus life}，但其他题目须重新审题。
\item \textbf{建议加作用：}\\
\en{The [组织] could organize [具体活动]. This would help [对象] [动词原形及其余部分].}\\
本题例：\en{The student union could organize sports events. This would help students stay healthy.} 可换活动与收益，但二者要有因果联系。
\item \textbf{邀请谁做什么：}\\
\en{It could invite [人] to [动词原形及其余部分].}\\
本题可以邀请毕业生分享工作经验。注意 \en{invite somebody to do}，与 \en{let somebody do} 不同。
\item \textbf{先了解需求，再采取行动：}\\
\en{The [组织] could conduct surveys to learn about [需求或困难]. Based on the results, it could [具体措施].}\\
适合征求意见、改善服务的任务；后一句必须说明调查结果怎样影响措施。
\end{enumerate}

\subsection{下一轮只练一个能力：建议加作用}
先不背整篇。选“运动”或“了解学生困难”中的一项，不看修改稿，用自己的话写两句：第一句写学生会能做什么，第二句写学生能得到什么。可以沿用 \en{could organize ... / This would help ...}，重点检查动词结构与活动、作用之间的联系。
